% Recovery provenance: reconstructed on 2026-11-20 from the complete
% authored mathematical text sent by higher_comparison2 in two messages.
% The historical 9343-character TeX byte stream was lost during the outage.
% All theorem/proof steps below were present in those independently audited chunks.
\Needspace{10\baselineskip}
\subsection{A.148. A relative parameter-Lipschitz factorization criterion}

\begin{saresult}{A.148}{Proved under stated hypotheses}{A pair of parameter-Lipschitz factors gives a PL representative}
Keep the real-field and fixed o-minimal structure conventions of A.147.
Let \(B\) be a compact closed PL manifold, and put
\(N=B\times I\), \(A_j=B\times\{j\}\), \(P=I^k\), \(k\ge1\).
Suppose \(r,s:N\times P\to N\) are continuous definable families of
homeomorphisms, each preserving \(A_0,A_1\) setwise and each uniformly
Lipschitz in the forward parameter variable. Suppose
\[
 r_p|_{A_0}=s_p|_{A_0}\quad(p\in P),\qquad
 r_p=s_p\quad(p\in\partial P).
 \tag{148.1}
\]
Then \(U_p=s_pr_p^{-1}\) is a based bottom-fixed cube and is based
definably homotopic, through bottom-fixed homeomorphisms, to a jointly
PL cube. Consequently the comparison image consists exactly of classes
admitting a representative with such a relative factorization.
Real semialgebraic \(T^6,\pi_2\) surjectivity is equivalent to every
based semialgebraic two-cube being based semialgebraically homotopic
to one admitting such a factorization with both forward parameter bounds.
\end{saresult}

\begin{proof}
All pair homotopies preserve (148.1), so quotients remain based and
bottom-fixed. Initially choose \(p_0\in\partial P\) and precompose both
factors by the fixed map \(r_{p_0}^{-1}\). The quotient, equality
conditions, and bounds are unchanged; both normalized factors are the
identity at \(p_0\). Connectedness of \(P\) makes them preserve every
component of \(N\).

\textbf{A single shared analytic and combinatorial construction.}
For a positive-dimensional connected component, double \(N\) to
\(X=N^+\cup N^-\) and reflect each factor. They preserve
\(N^\pm,A_0,A_1\) setwise. Choose the same marked \(Q\), the same
\(K_0\), the same fine Hudson \(K\) for both, and choose \(M\) using
\(\max(L_r,L_s)\). Thus \(f^j,R,\rho,T_v\) and the endpoint boxes
\(E_a\) are literally identical. For \(a=r,s\), solve
\[
 t_i^a=\rho_{a_{t^a}(x),v}(T_vp_i),\qquad
 a_p^v(x)=a_{t^a(x,p,v)}(x).
 \tag{148.2}
\]
A.147's residual lemma and common contraction bound give joint
continuous definability and exact inverse homeomorphism families.
On \(x\in A_0\) both systems coincide because \(r_t(x)=s_t(x)\) for
all \(t\). On \(p\in\partial P\), the endpoint rule forces \(t^a\)
into that same parameter face, where the full maps agree; the systems
again coincide. Hence equality persists on
\[
 E=(A_0\times P)\cup(X\times\partial P).
 \tag{148.3}
\]
At \(v=1\), the collections \(S_a\) of at most \(k\) support stars
and \(J_a\) depend only on the common \(K,M\). Both inverse formulas
give the same support identity outside \(J_a\). Apply Proposition 57
and A.147's finite strictification once, yielding one identical
presheaf \(\xi\) with the same finite parameter topology. Both factors
are \(\xi\)-prismatic and preserve all marks.

\textbf{One common target coloring and ambient change.}
Run the finite procedure of Mn\"ev's Lemma 8 with no auxiliary \(G_i\),
with \(Z=\varnothing\), and with the same constant marked \(Q\).
Its choices---closed PL approximation, finite order complex, split
marked rank construction, and the two mapping-cylinder colorings---
depend only on \(Q,\xi\), and the parameter topology, not on the input
map. Make them once. The inverse-image constancy proof verifies the
same terminal target coloring simultaneously for both factors.

Thus one base triangulation and one target collection
\(C\times\sigma\) have two finite definable source systems.
Their restrictions to \(E\) coincide: target cells in \(A_0\) have
identical preimages there, and on a parameter boundary face the full
maps agree. All mark incidences remain.

Include every spatial cell closure from both source systems, the
common target, and all \(Q\)-cells in one finite family. A.147's
compatible triangulation and marked Hauptvermutung give one
\(\phi\) preserving each \(Q\)-cell and making all these cells PL.
Its cellular isotopy from the identity preserves the marks.
Conjugate both factors by this same isotopy; equality on \(E\) remains,
so the quotients give a based bottom-fixed homotopy. Call the endpoint
factors \(V_r,V_s\). Both have finite PL prism systems with a common
target and identical restriction on the common subcomplex \(E\).

\textbf{Straighten the common restriction first.}
The common map \(V_E=V_r|_E=V_s|_E\) has one cell correspondence.
Construct one jointly PL \(G_E\) by total prism dimension, using
A.147's full-boundary PL extension. Extend \(G_E\) separately to
jointly PL \(G_r,G_s\) on the entire two source prism systems,
with their correspondences to the common target. Every new prism
meets \(E\) in a subcomplex of its full boundary, so \(G_E\) is
included in the prescribed data. Thus
\[
 G_r|_E=G_s|_E=G_E.
 \tag{148.4}
\]
Let \(K_r=G_r^{-1}V_r\) and \(K_s=G_s^{-1}V_s\).
Each preserves its source prisms; on \(E\) both equal
\(K_E=G_E^{-1}V_E\). First construct one cellwise definable
Alexander homotopy from \(K_E\) to \(\operatorname{id}_E\).
Extend this same prescribed homotopy separately to \(K_r\) and
\(K_s\) by total prism dimension. For a remaining \(C\times\sigma\),
let \(D=\sigma\times I\). On
\[
 (\partial C\times D)\cup(C\times\partial D)
\]
prescribe the already constructed face homotopies, initial \(K_a\),
final identity, and the common \(E\)-homotopy where applicable.
All data agree. A.147's formula
\[
 E_p=R_p A_{\|p\|_\infty}(Q_{p/\|p\|_\infty})
\]
fills the prism, with its \(2\|p\|_\infty\) displacement bound
proving continuity at the center. This extends the entire boundary
homotopy data; it does not require the initial \(V_E\) to be PL.

Write the resulting homotopies as \(K_r^w,K_s^w\), equal on \(E\).
Then \(R^w=G_rK_r^w\) and \(S^w=G_sK_s^w\) agree on \(E\), begin
at \(V_r,V_s\), end at \(G_r,G_s\), and preserve the halves and top
throughout. Their quotients \(S^w(R^w)^{-1}\) fix \(A_0\) pointwise
and are the identity on \(\partial P\).
The endpoint quotient is jointly PL: the total map
\(G_r:X\times P\to X\times P\) is a PL homeomorphism, so its inverse
and its composition with \(G_s\) are PL. Restrict to \(N^+\) and
concatenate the preceding pair homotopies.

\textbf{Components and converse.}
The initial normalization made the factors preserve every component
of \(N\). Apply the positive-dimensional proof componentwise and
paste the finitely many results. For a point component, the quotient
is an increasing interval homeomorphism, and the based convex
contraction of A.147 suffices. The empty case is immediate.
Conversely, any jointly PL based bottom-fixed cube \(G\) factors as
\(r=\operatorname{id}\), \(s=G\). Compactness gives both forward
bounds and (148.1). This proves the exact image characterization
and the stated \(T^6,\pi_2\) equivalence.
\end{proof}

\textbf{Scope and source distinction.}
The pair argument uses A.147's finite constructions. Its additional
point is that the proof of Mn\"ev, arXiv:0708.4039v3, Lemma 8,
printed pp.~67--68, makes choices from the coloring and configuration
presheaf, so common \(\xi\) gives one common target coloring.
Equality on \(E\) is retained by one initial straightening there and
relative extensions of its entire homotopy. This is not an arbitrary
common refinement of separately chosen target colorings.
The factors need not individually be based or bottom-fixed.
No spatial or inverse-parameter Lipschitz estimate, and no
Lipschitzness of the quotient, is assumed. Existence of the
factorization in every homotopy class is not asserted by this
criterion alone. The theorem as proved here has real ground field.

